% Complément au chapitre 1 de « Refaire l'impôt » : les élasticités estimées sur données françaises
% Compilation : pdflatex elasticites-france.tex (deux passes pour les renvois)
% Données de la figure : ../donnees/elasticites_france.csv ; figure : graphiques_cours.py, fonction c25
\documentclass[11pt,a4paper]{article}

\usepackage[T1]{fontenc}
\usepackage[utf8]{inputenc}
\usepackage{lmodern}
\usepackage[french]{babel}
\usepackage[a4paper,left=3.1cm,right=3.1cm,top=2.9cm,bottom=3.1cm]{geometry}
\usepackage{graphicx}
\usepackage[export]{adjustbox}
\usepackage{booktabs}
\usepackage{tabularx}
\usepackage{xcolor}
\usepackage{titlesec}
\usepackage[labelfont=bf,labelsep=period,font=small,singlelinecheck=false,
            justification=raggedright,skip=6pt]{caption}
\usepackage{fancyhdr}
\usepackage{enumitem}
\usepackage{microtype}
\usepackage{textcomp}
\providecommand{\euro}{\texteuro}
\usepackage[hidelinks]{hyperref}

\definecolor{encre}{HTML}{1A1A19}
\definecolor{bleu}{HTML}{003399}
\definecolor{vert}{HTML}{1F7A4D}
\definecolor{ocre}{HTML}{C79100}

\titleformat{\section}{\normalfont\large\bfseries\color{bleu}}{\thesection.}{0.6em}{}
\titleformat{\subsection}{\normalfont\normalsize\bfseries\color{bleu}}{}{0em}{}
\titlespacing*{\section}{0pt}{2.4ex plus 1ex minus .2ex}{1.4ex plus .2ex}
\titlespacing*{\subsection}{0pt}{2.2ex plus 1ex minus .2ex}{1ex plus .2ex}

\pagestyle{fancy}
\fancyhf{}
\renewcommand{\headrulewidth}{0.4pt}
\fancyhead[L]{\small\itshape Refaire l'impôt --- complément au chapitre 1}
\fancyhead[R]{\small\thepage}

\setlength{\parindent}{0pt}
\setlength{\parskip}{0.72em}
\linespread{1.06}
\renewcommand{\arraystretch}{1.22}

% Figure haute : légende au-dessus, image, puis sources et note de lecture en texte.
\newcommand{\figblocgrand}[3]{%
  \begin{figure}[htbp]\centering
  \caption{#2}\label{#3}
  \includegraphics[max width=\linewidth,max height=0.68\textheight]{../figures/nu/#1.pdf}
  \IfFileExists{../figures/nu/#1_note.tex}{\par\vspace{4pt}%
    \parbox{\linewidth}{\raggedright\scriptsize\color{encre}\setlength{\parskip}{2pt}%
    \input{../figures/nu/#1_note.tex}}}{}%
  \end{figure}}

\begin{document}
\addto\captionsfrench{\renewcommand{\figurename}{Figure}%
  \renewcommand{\tablename}{Tableau}}
\renewcommand{\figurename}{Figure}\renewcommand{\tablename}{Tableau}

\thispagestyle{empty}
{\footnotesize\color{ocre}\bfseries COMPLÉMENT AU CHAPITRE 1 DE «~REFAIRE L'IMPÔT~»\par}
\vspace{8pt}
{\raggedright\fontsize{21}{25}\selectfont\bfseries Les élasticités estimées\\ sur données françaises\par}
\vspace{8pt}
{\large Ce que disent les études publiées depuis 2010\par}
\vspace{10pt}
{\color{bleu}\hrule height 1.2pt}
\vspace{10pt}

Le chapitre~1 retient que le coût d'un impôt dépend de la sensibilité de son assiette au taux, et
que cette sensibilité tient en partie à ce que la loi permet. Douze travaux publiés depuis 2010 la
mesurent sur données françaises, mais ils ne mesurent pas tous la même chose. Six estiment de
combien un revenu déclaré réagit au taux net marginal qui le frappe ; la
figure~\ref{fig:elasticites} les rassemble, et ils dessinent une hiérarchie nette : le revenu
d'activité réagit peu, les revenus du capital davantage, les dividendes des dirigeants qui
contrôlent leur société bien plus encore. Les six autres mesurent la réaction d'autres grandeurs ---
patrimoine déclaré, lieu de résidence, montant dissimulé, dons, décision de travailler --- à
d'autres taux. Leurs chiffres ne se rangent pas sur la même échelle ; le tableau~\ref{tab:autres}
dit ce que chacun mesure. Pris un à un, ils montrent que les règles de déclaration et de contrôle
pèsent autant que le taux.

\section{Ce que l'on mesure}

L'élasticité rapporte la variation en pourcentage du revenu déclaré à la variation en pourcentage
de ce qui reste au contribuable sur l'euro marginal, le taux net $1-\tau$. Elle donne directement
le coût de l'impôt : la perte sèche par euro prélevé en relevant un taux vaut
$e\,\tau/(1-\tau-e\,\tau)$ (Saez, Slemrod et Giertz, 2012), et le taux supérieur qui maximise les
recettes vaut $1/(1+a\,e)$, où $a$ mesure l'épaisseur de la queue haute de la distribution.

Le revenu déclaré peut bouger pour trois raisons. Le contribuable travaille, investit ou épargne
moins : c'est la réponse réelle. Il déclare autrement le même revenu --- plus tard, sous une autre
forme, dans une société, ou pas du tout. Il sort de l'assiette : il cesse de travailler, ou il quitte
le pays. Les études mesurent la somme ; seules certaines parviennent à séparer les composantes. C'est
pourtant la séparation qui compte pour la politique fiscale, puisque la deuxième réponse dépend des
règles et non des préférences.

Trois méthodes coexistent. La plus répandue régresse la variation du revenu d'une année sur l'autre
sur la variation du taux net, en instrumentant celle-ci par la variation qu'aurait produite la
réforme seule si le revenu n'avait pas bougé. Elle doit neutraliser le retour à la moyenne : un
revenu exceptionnellement haut une année a tendance à baisser l'année suivante, impôt ou pas, et
l'estimateur prend ce mouvement pour une réaction au taux si on ne le contrôle pas (Gruber et Saez,
2002 ; Weber, 2014). La deuxième compte les contribuables qui s'accumulent juste sous un seuil où le
taux augmente : l'identification est propre mais locale, et l'inattention ou les coûts d'ajustement
la tirent vers le bas. La troisième compare, avant et après une réforme, un groupe touché et un
groupe épargné, puis divise l'écart d'évolution du revenu par l'écart d'évolution du taux net.
Presque toutes les estimations françaises sont de court terme : elles portent sur des variations
d'une année sur l'autre, et ne disent rien de ce qui se passe en cinq ou dix ans.

\figblocgrand{c25_elasticites_france}{La réaction des revenus déclarés au taux net marginal qui les
frappe, études françaises depuis 2010}{fig:elasticites}

\section{Revenu imposable et revenus d'activité}

\subsection{Cabannes, Houdré et Landais (2014) : presque rien en moyenne, 0,31 en haut}

La question est celle de la réaction du revenu imposable des foyers aux baisses d'impôt sur le
revenu de la fin des années 1990 et du début des années 2000. Les données sont les échantillons
annuels de 500\,000 déclarations de la DGFiP, qui contiennent le revenu de l'année et celui de la
précédente, pour les revenus 1997 à 2004. Deux réformes fournissent la variation : la baisse du
plafond du quotient familial entre 1997 et 1998, qui a changé le taux marginal de foyers de même
revenu selon le nombre de leurs enfants, et les baisses successives des taux du barème entre 2000 et
2003. L'instrument est la variation de taux qu'aurait produite la législation nouvelle sur le revenu
de l'année précédente.

L'élasticité compensée vaut 0,02 sur l'ensemble des foyers, premier décile exclu : à peu près rien.
Elle n'est pas significative entre le premier et le neuvième décile, et vaut 0,31 pour le décile
supérieur (0,50 avec une autre pondération et un effet revenu). Les auteurs, qui ne publient que
des seuils de significativité et pas d'écart-type, encadrent cette valeur centrale par deux
scénarios à 0,15 et 0,45 ; ils en tirent, à titre indicatif, un taux maximisant les recettes de
61~\% sur l'assiette de 2004, entre 76 et 51~\% selon le scénario. Deux limites tiennent aux
données : en haut de la distribution, les baisses de taux ont été presque identiques pour tous les
foyers, ce qui fragilise l'identification ; et deux années consécutives ne disent rien des réponses
de long terme.

\subsection{Lehmann, Marical et Rioux (2013) : l'impôt sur le revenu, pas les cotisations}

La question est différente : un salaire réagit-il de la même façon à l'impôt sur le revenu et aux
cotisations sociales ? La plupart des modèles du marché du travail le prédisent, puisque seul y
compte l'écart entre le coût du travail et ce que touche le salarié. Les enquêtes Revenus fiscaux de
2003 à 2006, qui apparient l'enquête Emploi et les déclarations d'impôt, suivent environ
12\,500~salariés du privé sur deux années consécutives. Deux réformes qui ne touchaient que les
salariés gagnant moins de deux fois le salaire minimum fournissent la variation : la forte hausse de
la prime pour l'emploi, dont le montant maximal a presque doublé entre 2003 et 2006, et la
convergence des deux barèmes d'allègements de cotisations patronales, celui des entreprises restées
à 39~heures et celui des entreprises passées à 35~heures.

Le salaire réagit au taux marginal de l'impôt sur le revenu avec une élasticité d'environ 0,2
(0,225 avec tous les contrôles, écart-type 0,085, dans la version de travail), et pas du tout au
taux marginal des cotisations. La réponse vient des femmes (0,88) ; parmi les salariés modestes
employés toute l'année les deux années, elle tombe à 0,03. Elle passe donc par le nombre de mois
travaillés plutôt que par les heures. Les auteurs expliquent l'asymétrie par la rigidité de court
terme du salaire horaire, négocié à intervalles espacés : la prime pour l'emploi joue sur la décision
de travailler, alors que les cotisations pèsent sur un salaire qui ne se renégocie pas chaque année.
Une baisse des cotisations patronales réduit ainsi, à court terme, le coût du travail sans modifier
le salaire. L'élasticité est identifiée sur des salariés modestes ; elle ne dit rien du haut de la
distribution.

\subsection{Sicsic (2022) : trois fois plus de réaction à l'impôt qu'aux prestations}

Sicsic prolonge la comparaison aux prestations. Avec les enquêtes Revenus fiscaux et sociaux de 2006
à 2015, soit environ 90\,000~personnes suivies deux années de suite, et les nombreuses réformes de la
période --- barème de 2007, tranche à 45~\% en 2013, gel du barème de 2011 à 2013, décote, plafond
du quotient familial, création du RSA activité en 2009, modulation des allocations familiales ---,
il estime l'élasticité compensée du revenu d'activité au taux net marginal de chaque dispositif. Elle
vaut 0,26 pour l'impôt sur le revenu (écart-type 0,018), 0,09 pour le RSA activité et zéro pour les
prestations familiales ; 0,09 pour l'ensemble des transferts. Elle est plus forte en haut de la
distribution (0,51 pour les 10~\% de revenus d'activité les plus élevés), chez les personnes seules
sans enfant (0,53) et chez les indépendants (0,43, très imprécise), plus faible chez les couples avec
enfants (0,17).

Un même taux marginal produit donc une réaction trois fois plus forte quand il vient de l'impôt sur
le revenu que quand il vient du retrait d'une prestation. L'auteur y voit la saillance de l'impôt,
mieux perçu que les règles de calcul des prestations. Les revenus d'activité, déclarés par
l'employeur et sans déduction, se prêtent mal au déplacement d'assiette : ces élasticités sont plus
proches d'une réponse réelle que celles du revenu imposable total. Elles sont de court terme, et
l'auteur les juge plutôt sous-estimées, parce que les réponses prennent du temps.

\subsection{Lardeux (2018) : une réaction déclarative à un barème mal compris}

Lardeux observe un cas où le revenu déclaré se choisit presque librement : les célibataires qui
reçoivent de leur famille une obligation alimentaire et en déclarent eux-mêmes le montant. Leurs
déclarations s'accumulent au vrai point d'entrée dans l'impôt, le seuil de mise en recouvrement,
mais aussi à un autre seuil du barème que beaucoup prennent pour lui. En combinant les deux
accumulations dans un modèle de perception erronée, il estime une élasticité d'environ 0,7 au taux
marginal perçu et une attention d'environ 75~\% au vrai barème ; environ 80~\% des déclarants
concernés n'ajustent pas leur déclaration. La réponse est purement déclarative et ne concerne qu'une
petite population ; elle montre surtout qu'une partie des contribuables optimise sur un barème
qu'elle comprend mal. Ce 0,7 ne se compare pas aux chiffres de la figure~\ref{fig:elasticites} :
c'est l'élasticité des seuls déclarants qui ajustent, corrigée des frictions, au taux qu'ils croient
payer, alors que les autres études mesurent une réaction moyenne sur tous les contribuables, y
compris ceux qui ne bougent pas.

\section{Revenus du capital}

\subsection{Lefebvre, Lehmann et Sicsic (2025) : 0,77, et un effet croisé qui compte plus}

La question est celle du taux qui maximise les recettes sur les revenus du capital. Les auteurs
montrent qu'il dépend de deux élasticités : celle des revenus du capital à leur propre taux net, et
celle des revenus du travail à ce même taux. La seconde paraît secondaire ; elle ne l'est pas, parce
que les revenus du travail sont bien plus gros que ceux du capital, si bien qu'une petite réaction
des salaires à la fiscalité du capital pèse lourd dans la recette totale.

Les données sont les déclarations exhaustives de l'impôt sur le revenu de 2008 à 2017, suivies foyer
par foyer : 2,3~millions de foyers stables, dont le revenu fiscal de référence dépasse 30\,000~\euro{}
et qui déclarent chaque année des revenus du travail et du capital. L'effet est identifié par les
variations du taux du prélèvement forfaitaire libératoire, que les foyers pouvaient choisir de 2008 à
2012 pour leurs dividendes et leurs intérêts, et par la suppression de cette option en 2013, qui a
soumis ces revenus au barème.

L'élasticité directe vaut 0,767 (écart-type 0,011) dans la spécification préférée. Les élasticités
croisées des revenus du travail sont positives, de 0,036 à 0,083 selon la tranche : quand l'impôt sur
le capital baisse, les foyers déclarent aussi un peu plus de revenus du travail, ce qui exclut un
simple transfert des salaires vers les dividendes. Sans l'effet croisé, le taux de Laffer serait de
56,6~\% ; avec lui, de 42,9~\%, avant même de compter les cotisations sociales.

La précision affichée trompe sur l'incertitude réelle. Avec les instruments de Weber, qui corrigent
mieux le retour à la moyenne, l'élasticité directe tombe à 0,51, voire 0,22 selon les contrôles de
revenu ; d'autres variantes la portent à 1,17. Le rapport que les mêmes auteurs ont remis à France
Stratégie en 2020, avec Zanoutene, trouvait 0,67 en spécification de base, 0,30 en différences de
deux ans et 1,17 avec les instruments de Weber. L'ordre de grandeur, entre 0,3 et 1, est solide ; le
chiffre exact ne l'est pas.

\subsection{Bach et al. (2024) : le dirigeant laisse le bénéfice dans la société}

Bach, Bozio, Guillouzouic, Leroy et Malgouyres cherchent ce que recouvrent ces réactions. Ils
apparient les déclarations d'impôt sur le revenu et d'ISF aux liasses fiscales des sociétés dont les
foyers détiennent au moins 10~\%, ce qui permet de suivre l'argent de l'entreprise au ménage. Deux
réformes de sens opposé fournissent la variation : la suppression du prélèvement libératoire en 2013,
qui a relevé le taux supérieur sur les dividendes de 36,5 à 40,2~\%, et le prélèvement forfaitaire
unique de 2018, qui l'a ramené à 30~\%. Parmi les assujettis à l'ISF, ils comparent les foyers que
leurs revenus réguliers placent dans les tranches supérieures, touchés par les réformes, aux autres.

Les dividendes des foyers qui contrôlent une société baissent de 20~\% en 2013 et montent de 20~\% en
2018. Rapporté à la variation du taux net, cela donne une élasticité de 2,5 pour la baisse de 2018,
que les auteurs jugent la mieux mesurée, et de 3,2 pour la hausse de 2013. Celle des foyers sans
société est nulle en 2018. La différence tient à la capacité d'agir plutôt qu'aux préférences : le
dirigeant décide de la distribution. Du côté des sociétés, l'investissement ne bouge dans aucun des
deux cas. En 2013, la baisse des dividendes est passée par une hausse de l'épargne des sociétés,
puis par une baisse du bénéfice imposable des plus petites ; en 2018, la baisse de l'épargne des
sociétés explique à elle seule la hausse des dividendes. Les dirigeants traitent leur société comme
un placement peu taxé, où le bénéfice peut attendre que le dividende soit moins imposé. Ni un
déplacement vers d'autres revenus du capital, ni une substitution entre salaire et dividende
n'expliquent plus qu'une petite part de la réaction. Les élasticités sont publiées sans écart-type,
et la détention de sociétés est mal mesurée pour 2013.

Une élasticité de 2,5, prise telle quelle, donnerait un taux maximisant les recettes de 29~\% sur les
dividendes. Elle mesure pourtant surtout un report : le bénéfice laissé dans la société reste
imposable plus tard. L'élasticité de long terme est probablement plus faible, mais personne ne l'a
mesurée.

\subsection{Ce qu'il ne faut plus citer, et ce que conclut France Stratégie}

Boissel et Matray (\emph{American Economic Review}, 2022) avaient conclu, à partir de la réforme de
2013, que les entreprises dont les dividendes étaient plus taxés avaient réduit leurs distributions
et investi davantage. Un commentaire de Bach, Bozio, Guillouzouic et Malgouyres (2023) a relevé une
altération du code produisant le graphique d'étude d'événement et des tendances différentes, avant la
réforme, entre entreprises traitées et témoins ; les auteurs ont retiré l'article en juillet 2023.
Lefebvre, Lehmann et Sicsic le citent encore comme point de comparaison ; leurs propres résultats
n'en dépendent pas.

Le comité d'évaluation des réformes de la fiscalité du capital conclut en 2023 que le prélèvement
forfaitaire unique a causé la majeure partie de la hausse des dividendes de 2018-2019, sans pouvoir
en chiffrer l'ordre de grandeur macroéconomique, et ne détecte d'effet ni sur l'investissement ni sur
les salaires.

\section{Autres assiettes et autres marges}

Les études de cette partie ne mesurent pas la réaction d'un revenu déclaré au taux net qui le
frappe. Elles mesurent celle d'autres grandeurs, à d'autres taux, et leurs chiffres ne se comparent
ni entre eux ni avec ceux de la figure~\ref{fig:elasticites} ; le tableau~\ref{tab:autres} dit ce
que chacun mesure. Ce qui les rapproche est une leçon que chacune établit à l'intérieur de son propre
cadre : changer les règles de déclaration ou de contrôle déplace autant, ou plus, que changer le
taux.

\begin{table}[htbp]
\caption{Les élasticités qui ne se placent pas sur la figure~\ref{fig:elasticites}}
\label{tab:autres}
\footnotesize
\renewcommand{\arraystretch}{1.25}
\begin{tabularx}{\linewidth}{@{}>{\raggedright\arraybackslash}p{3.1cm}
  >{\raggedright\arraybackslash}X>{\raggedright\arraybackslash}X>{\raggedleft\arraybackslash}p{1.9cm}@{}}
\toprule
\textbf{Étude} & \textbf{Ce qui réagit} & \textbf{À quoi} & \textbf{Estimation}\\
\midrule
Lardeux (2018) & revenu déclaré des seuls contribuables qui ajustent & taux marginal qu'ils croient
payer & 0,7\\
Lefebvre, Lehmann et Sicsic (2025) & revenus du travail & taux net sur les revenus du capital
(élasticité croisée) & 0,04 à 0,08\\
Garbinti \emph{et al.} (2026) & croissance annuelle du patrimoine déclaré & taux net de l'ISF, un
moins le taux annuel & 0,13\\
Focus du CAE (2025) & nombre de résidents parmi les 1~\% les plus aisés & revenu net d'impôt & 0,02
à 0,18\\
Aghion \emph{et al.} (2023) & montant dissimulé par les indépendants & taux d'imposition lui-même &
1,1 à 1,2\\
Fack et Landais (2016) & dons déclarés & prix du don, un moins le taux marginal d'imposition & $-1{,}86$
puis $-0{,}36$\\
Carbonnier (2021) & décision de travailler des femmes mariées & part gardée du revenu de la reprise
d'emploi & 1 à 2\\
\bottomrule
\end{tabularx}
\end{table}

\subsection{Garbinti et al. (2026) : l'ISF réagit à l'information plus qu'au taux}

Garbinti, Goupille-Lebret, Muñoz, Stantcheva et Zucman étudient l'impôt sur la fortune avec
l'ensemble des déclarations de 2006 à 2017, appariées à l'impôt sur le revenu. Aux seuils de tranche,
où seul le taux change, les contribuables bougent à peine : l'élasticité de la croissance du
patrimoine déclaré au taux net de l'impôt vaut 0,13 (écart-type 0,04) au deuxième seuil et n'est pas
significative au troisième. Ils réagissent en revanche à l'information demandée. Depuis 2011, les
foyers dont le patrimoine était inférieur à un seuil --- 3~millions d'euros, puis 2,57~millions à
partir de 2013 --- pouvaient déclarer un montant global sans détailler leurs actifs. Après
l'abaissement du seuil, les foyers proches ont ralenti la croissance déclarée de leur patrimoine pour
rester dessous : de 0,28~point par an en moyenne, dix fois la réaction aux seuils de tranche. Pour
obtenir le même effet par le taux, il aurait fallu porter le taux marginal de 0,7 à 2,84~\%. Les
foyers restés sous le seuil ont sous-déclaré 18,4~\% de l'impôt qu'ils auraient payé. Les indices
--- montants ronds, immobilier plus touché que les actifs financiers, revenus déclarés par des tiers
qui ne bougent pas, sur un test que les auteurs jugent peu puissant --- désignent la sous-déclaration
plutôt qu'une épargne moindre. La version lue est la révision de mai 2026 du document de travail.

\subsection{Aghion et al. (2023) : les régimes simplifiés facilitent la fraude}

Aghion, Akcigit, Gravoueille, Lequien et Stantcheva observent, sur l'ensemble des déclarations de
2006 à 2015, une forte accumulation de travailleurs indépendants juste sous les plafonds de recettes
des régimes simplifié, créé en 1999, et super-simplifié, créé en 2009. Plusieurs indices désignent la
fraude : déclarants qui restent anormalement dans la même tranche de recettes d'une année sur l'autre,
montants ronds, recettes réparties dans le couple. Leur modèle, qui suppose nulle la réponse réelle
parce que l'accumulation ne s'explique pas par elle, l'attribue au goût de la simplicité ---
l'équivalent de 1\,100 à 1\,600~\euro{} par an pour le régime simplifié, 6\,800 à 7\,200~\euro{} pour
le super-simplifié --- et à une élasticité du montant dissimulé au taux d'imposition d'environ 1,1.

\subsection{Fack et Landais (2016) : la même déduction, cinq fois moins élastique sous contrôle}

En 1983, l'administration a exigé que les reçus des dons soient joints à la déclaration, alors qu'il
suffisait auparavant de les conserver. Les dons déclarés ont chuté, sans rupture dans les dons reçus
par la Fondation de France : la baisse venait donc surtout de dons déclarés qui n'avaient pas été
faits. L'élasticité-prix des dons déclarés est passée de $-1{,}86$ à $-0{,}36$. Les données ---
échantillons de déclarations de 1975 à 1988 --- sont anciennes, mais le résultat est le plus net de
cette revue : la même déduction, au même taux, devient cinq fois moins élastique quand
l'administration vérifie.

\subsection{Le Focus du CAE (2025) : des départs sensibles à l'impôt, mais rares}

Bach, Bozio, Grimprel, Guillouzouic, Landais et Malgouyres mesurent les départs à l'étranger des
foyers les plus aisés, avec les déclarations d'impôt sur le revenu et d'ISF ou d'IFI de 2010 à 2019.
Chaque année, 0,2~\% des 1~\% de foyers aux revenus du capital les plus élevés quittent la France. La
hausse d'impôt de 2012-2013 a accru les départs du groupe le plus touché de 0,04 à 0,09~point ; les
baisses de 2017-2018 les ont réduits de 0 à 0,02~point et ont accru les retours de 0,01 à 0,04~point.
Rapportées à la variation du revenu net d'impôt, ces réactions donnent une élasticité de long terme
du nombre de résidents comprise entre 0,02 et 0,18. L'effet sur la valeur ajoutée des entreprises
détenues par les partants est très faible, et les auteurs jugent que les réponses de ceux qui restent
--- optimisation, épargne, fraude --- pèsent bien davantage. Le Focus ne publie pas d'écart-type.

\subsection{La participation au marché du travail}

Carbonnier (2021) exploite les discontinuités du barème que crée l'imposition commune des couples
mariés pour estimer l'élasticité de la participation des femmes mariées à leur taux de rétention à
la reprise d'emploi : entre 1 et 2, plus forte en bas de la distribution, en forme de U le long des
revenus. Elle porte sur la décision de travailler et non sur le revenu déclaré ; elle n'est connue
ici que par le résumé publié, le texte n'étant pas accessible depuis l'environnement de travail. Les
travaux sur le RSA recensés par Sicsic (2023) trouvent des effets faibles ou nuls sur l'emploi
(Briard et Sautory, 2012 ; Bargain et Vicard, 2014), sauf pour les mères isolées de jeunes enfants
(Simonnet et Danzin, 2014) ; ils ne donnent pas d'élasticité.

\section{Ce qu'on peut en conclure}

Parmi les estimations comparables, les revenus du capital réagissent environ trois fois plus que le
revenu d'activité, et les dividendes de ceux qui contrôlent leur société dix fois plus ; ces
rapports sont approximatifs, puisque les populations et les horizons diffèrent. Cette hiérarchie
suit la facilité de changer la
forme, la date ou le lieu d'un revenu sans changer d'activité. Un salaire déclaré par
l'employeur ne se déplace pas ; un dividende que le dirigeant peut laisser dans sa société se
déplace d'une année à l'autre ; un patrimoine déclaré en un seul chiffre, ou un don qu'on n'a pas à
justifier, se déplace dans la déclaration elle-même.

Les études françaises donnent ainsi à la phrase de Slemrod et Kopczuk (2002) --- l'élasticité est en
partie un choix de politique publique --- des preuves directes, chacune obtenue en comparant deux
situations dans une même étude. Supprimer le détail de la déclaration d'ISF a produit dix fois
l'effet des seuils de tranche ; exiger les reçus des dons a divisé
l'élasticité par cinq ; les régimes simplifiés des indépendants s'accompagnent de fraude. La
saillance compte aussi : un même taux marginal produit plus de réaction quand il vient de l'impôt sur
le revenu que des cotisations ou des prestations.

Trois manques limitent ce qu'on peut en tirer. Toutes les élasticités sont de court terme, mesurées
d'une année sur l'autre ; personne n'a estimé en France la réaction à cinq ou dix ans. Aucune étude
récente n'estime la réaction du revenu d'activité du 1~\% le plus aisé, alors que c'est elle qui
décide du taux supérieur. Et la part réelle des réactions n'est séparée de la part déclarative que
dans quelques cas, toujours par des indices plutôt que par une mesure.

\vspace{1.2em}
{\color{bleu}\hrule height 0.8pt}
\vspace{0.8em}

{\small
\textbf{Études recensées.} Laurent Bach, Antoine Bozio, Arthur Guillouzouic, Claire Leroy et Clément
Malgouyres, «~Follow the money! Why dividends overreact to flat-tax reforms~», document de travail,
septembre 2024 \textbullet{} Laurent Bach, Antoine Bozio, Nicolas Grimprel, Arthur Guillouzouic,
Camille Landais et Clément Malgouyres, «~Fiscalité du capital : quels sont les effets de l'exil
fiscal sur l'économie ?~», \emph{Focus du CAE} n°~118, juillet 2025 \textbullet{} Pierre-Yves
Cabannes, Cédric Houdré et Camille Landais, «~Comment le revenu imposable des ménages réagit-il à sa
taxation ? Une estimation sur la période 1997-2004~», \emph{Économie et Statistique} n°~467-468,
2014 \textbullet{} Clément Carbonnier, «~Imposition jointe des revenus et emploi des femmes mariées :
estimation à partir du cas français~», \emph{Revue économique} 72(2), 2021 \textbullet{} Gabrielle
Fack et Camille Landais, «~The effect of tax enforcement on tax elasticities: Evidence from
charitable contributions in France~», \emph{Journal of Public Economics} 133, 2016 (lu : document de
travail, novembre 2013) \textbullet{} Bertrand Garbinti, Jonathan Goupille-Lebret, Mathilde Muñoz,
Stefanie Stantcheva et Gabriel Zucman, «~Tax Design, Information, and Elasticities: Evidence From
the French Wealth Tax~», NBER \emph{Working Paper} 31333, révisé en mai 2026 \textbullet{} Philippe
Aghion, Ufuk Akcigit, Maxime Gravoueille, Matthieu Lequien et Stefanie Stantcheva, «~Tax Simplicity
or Simplicity of Evasion? Evidence from Self-Employment Taxes in France~», document de travail, mai
2023 \textbullet{} Raphaël Lardeux, «~Who Understands The French Income Tax? Bunching Where Tax
Liabilities Start~», document de travail Insee G2018/04, 2018 \textbullet{} Marie-Noëlle Lefebvre,
Etienne Lehmann et Michaël Sicsic, «~Estimating the Laffer tax rate on capital income: cross-base
responses matter!~», \emph{Scandinavian Journal of Economics}, 2025 (lu : IZA DP 16112, 2023)
\textbullet{} Marie-Noëlle Lefebvre, Etienne Lehmann, Michaël Sicsic et Elsa Zanoutene, «~Évaluation
de la mise au barème des revenus du capital~», rapport pour France Stratégie, septembre 2020
\textbullet{} Etienne Lehmann, François Marical et Laurence Rioux, «~Labor income responds
differently to income-tax and payroll-tax reforms~», \emph{Journal of Public Economics} 99, 2013
(lu : IZA DP 6108, 2011) \textbullet{} Michaël Sicsic, «~Does labour income react more to income tax
or means-tested benefits reforms?~», \emph{Fiscal Studies} 43(3), 2022 (lu : TEPP WP 2020-03).

\textbf{Autres travaux cités.} Laurent Bach, Antoine Bozio, Arthur Guillouzouic et Clément
Malgouyres, «~Dividend Taxes and the Allocation of Capital: Comment~», \emph{American Economic
Review} 113(7), 2023 \textbullet{} Charles Boissel et Adrien Matray, «~Dividend Taxes and the
Allocation of Capital~», \emph{American Economic Review} 112(9), 2022, retiré par les auteurs,
\emph{American Economic Review} 113(7), 2023 \textbullet{} Jonathan Gruber et Emmanuel Saez, «~The
elasticity of taxable income: evidence and implications~», \emph{Journal of Public Economics}, 2002
\textbullet{} Emmanuel Saez, Joel Slemrod et Seth Giertz, «~The Elasticity of Taxable Income with
Respect to Marginal Tax Rates: A Critical Review~», \emph{Journal of Economic Literature}, 2012
\textbullet{} Michaël Sicsic, «~L'élasticité de l'offre de travail et des revenus dans la
littérature~», \emph{Revue de l'OFCE} n°~183, 2023 \textbullet{} Joel Slemrod et Wojciech Kopczuk,
«~The optimal elasticity of taxable income~», \emph{Journal of Public Economics}, 2002
\textbullet{} France Stratégie, Comité d'évaluation des réformes de la fiscalité du capital,
quatrième rapport, avis du comité, octobre 2023 \textbullet{} Caroline Weber, «~Toward obtaining a
consistent estimate of the elasticity of taxable income using difference-in-differences~»,
\emph{Journal of Public Economics}, 2014.

\textbf{Méthode.} Chaque chiffre a été relu dans le document cité ; le fichier
\texttt{elasticites\_france.csv}, dans le dossier \texttt{donnees}, donne pour chacun la page et le
tableau. Quand la version publiée n'était pas accessible, la version
de travail lue est indiquée ; ses chiffres peuvent différer de ceux de l'article.}

\end{document}
